\documentclass[10pt,a4paper]{article}

\usepackage[margin=1in]{geometry}
\usepackage{amsmath,amssymb,amsthm}
\usepackage{hyperref}
\usepackage{booktabs}
\usepackage{algorithm}
\usepackage{algpseudocode}
\usepackage{float}
\usepackage{xcolor}
\usepackage{url}

\hypersetup{colorlinks=true, linkcolor=blue!50!black, citecolor=blue!50!black, urlcolor=blue!50!black}

\newtheorem{proposition}{Proposition}
\newtheorem{lemma}{Lemma}
\theoremstyle{remark}
\newtheorem{remark}{Remark}

\newcommand{\FF}{\mathbb{F}}
\newcommand{\ZZ}{\mathbb{Z}}
\newcommand{\Fb}{\mathcal{F}}
\newcommand{\Sem}[1]{S_{#1}}
\newcommand{\abs}[1]{\lvert #1 \rvert}
\newcommand{\code}[1]{\texttt{#1}}

\title{Point Decomposition by Root Finding:\\
       Replacing Gr\"obner-Basis Solvers\\
       in Elliptic-Curve Index Calculus\\[4pt]
       {\large with Measurements on Prime and Binary Curves up to NIST P-521}}
\author{Anonymous}
\date{}

\begin{document}
\maketitle

\begin{abstract}
The decomposition oracle at the heart of elliptic-curve index calculus
asks, for a target point $R$ and a factor base $\Fb$, whether
$R = P_1 + \cdots + P_m$ with every $P_i \in \Fb$. The textbook answer
reduces the question to the roots of Semaev's summation polynomial and
hands the resulting polynomial system to a Gr\"obner-basis engine
(F4/F5). We observe that for the summation polynomial $\Sem{3}$ the
roots are already known in closed form---they are the $x$-coordinates of
$R \mp P_i$---so the polynomial solve can be replaced by direct
evaluation of the group law followed by a table lookup, and for binary
curves by a meet-in-the-middle table of pair sums. Both replacements are
exact and complete for the decomposition sizes they address, and both
admit a fail-closed verification in which every reported decomposition
is re-added in the group. We implement them, make root finding the
default decomposition strategy in an existing index-calculus codebase,
and measure against the same codebase's Gr\"obner solvers on identical
targets: $383$--$1535\times$ per decomposition on Koblitz curves over
$\FF_{2^9}$ and $\FF_{2^{13}}$ (versus matrix-F4), $3\times 10^4$--$5\times 10^4\times$
on small prime curves, and $3021$--$5802\times$ on the five NIST prime
curves P-192, P-224, P-256, P-384 and P-521 with their standard
parameters (versus Buchberger on the same decomposition ideal, in the
same field). A dedicated P-256 field implementation gives a further
$11$--$12\times$ over generic arbitrary-precision arithmetic.

We are explicit about what this does and does not change. The result is
an improvement in the \emph{constant} of a subroutine. It leaves the
asymptotic cost of index calculus unchanged, does not alter the
exponential size of the Frobenius-invariant factor base on Koblitz
curves, does not make prime-field index calculus competitive with
Pollard rho, and makes no claim about the security of any deployed
curve. All instances are public synthetic or standard-parameter
known-answer instances; all code and raw measurements are released.
\end{abstract}

\section{Introduction}
\label{sec:intro}

Index calculus on elliptic curves, in the form proposed by
Semaev~\cite{semaev2004} and developed by Gaudry~\cite{gaudry2009} and
Diem~\cite{diem2011}, has two expensive stages. The first collects
\emph{relations}: random multiples $R = [r]P + [s]Q$ are tested for
membership in the $m$-fold sumset of a factor base $\Fb$, and each hit is
a linear relation among the unknown discrete logarithms of $\Fb$. The
second solves the resulting sparse linear system. The first stage is
dominated by the \emph{decomposition oracle}: deciding whether a given $R$
is a sum of $m$ factor-base points and, if so, finding the summands.

The standard route to this oracle is algebraic. Semaev's summation
polynomial $\Sem{m+1}(x_1, \ldots, x_m, x_R)$ vanishes exactly when
there exist signs such that $\pm P_1 \pm \cdots \pm P_m = R$ for points
with the given $x$-coordinates. One restricts the $x_i$ to the factor
base---either by adjoining vanishing polynomials $\prod_{F\in\Fb}(x_i -
x_F)$ over a prime field, or by Weil descent onto a subspace over an
extension field---and solves the resulting zero-dimensional polynomial
system. The solver of choice is a Gr\"obner-basis algorithm:
Faug\`ere's F4~\cite{faugere1999} or F5~\cite{faugere2002}, often in a
matrix formulation tuned to the system's shape. Substantial effort has
gone into making this solve tractable, through symmetrization, the
exploitation of Frobenius and sign symmetries, and careful degree
analysis~\cite{fppr2012,petitquisquater2012,galbraithgebregiyorgis2014}.

This paper starts from a simple observation that is well known but, in
our experience, under-exploited in implementations: for $m = 2$ the
polynomial system does not need to be solved, because its variety is
already parametrized. The roots of $\Sem{3}(x_R, x_i, X)$ in $X$ are
exactly the $x$-coordinates of $R + F_i$ and $R - F_i$, where $F_i$ is
any factor-base point with $x$-coordinate $x_i$. One field inversion
(amortized across the factor base by Montgomery's trick) and a hash
lookup per factor-base element therefore decide the $2$-decomposition of
$R$ exactly. Over binary fields the same idea takes the form of a
meet-in-the-middle table of pair sums, which decides $2$-, $3$- and
$4$-decompositions in $1$, $\abs{\Fb}$ and $\abs{\Fb}^2$ lookups
respectively. Neither method approximates; both are complete for the
sizes they address, and both report a witness that can be re-added in
the group.

\paragraph{Contributions.}
\begin{enumerate}
\item We state the two root-finding decomposition procedures precisely
  (\S\ref{sec:algorithms}), prove their exactness
  (Proposition~\ref{prop:s3roots}), and give their cost in field
  operations and lookups.
\item We implement both in an existing Rust index-calculus codebase that
  already contained matrix-F4 and Buchberger decomposition solvers, make
  root finding the default strategy, and keep the Gr\"obner solvers as
  selectable baselines (\S\ref{sec:impl}).
\item We define a fail-closed verification protocol under which every
  reported decomposition is re-added in the group, cross-checked against
  the summation polynomial, and---where a Gr\"obner basis is also
  computed---checked to lie in the variety of that basis
  (\S\ref{sec:verification}).
\item We measure root finding against the Gr\"obner solvers on identical
  targets (\S\ref{sec:experiments}): on binary Koblitz curves, on small
  prime curves, on prime fields from $40$ to $83$ bits, and on the five
  NIST prime curves P-192 through P-521 with their standard parameters,
  including a per-curve specialized field for P-256.
\item We state precisely what the measurements do and do not imply
  (\S\ref{sec:boundary}). In particular, the asymptotics of index
  calculus are unchanged, the exponential factor-base wall on Koblitz
  curves is untouched, and no security claim about any deployed curve
  follows.
\end{enumerate}

\section{Background}
\label{sec:background}

\subsection{Summation polynomials and the decomposition problem}

Let $E / \FF_q$ be an elliptic curve. Semaev~\cite{semaev2004} defined
the summation polynomials $\Sem{m} \in \FF_q[x_1,\ldots,x_m]$ by the
property that for $x_1,\ldots,x_m \in \overline{\FF}_q$,
\[
  \Sem{m}(x_1,\ldots,x_m) = 0
  \iff
  \exists\, P_1,\ldots,P_m \in E(\overline{\FF}_q):\;
  x(P_i) = x_i,\;\; P_1 + \cdots + P_m = \mathcal{O}.
\]
For short Weierstrass $y^2 = x^3 + ax + b$ the third polynomial is
\begin{equation}
\label{eq:s3}
\Sem{3}(x_1,x_2,x_3) =
  (x_1 - x_2)^2 x_3^2
  - 2\big[(x_1+x_2)(x_1 x_2 + a) + 2b\big] x_3
  + \big[(x_1 x_2 - a)^2 - 4b(x_1+x_2)\big],
\end{equation}
and higher $\Sem{m}$ are obtained by resultants, with
$\deg_{x_i} \Sem{m} = 2^{m-2}$.

Given a factor base $\Fb \subset E(\FF_q)$ and a target $R$, the
\emph{$m$-decomposition problem} is to decide whether there exist
$F_1,\ldots,F_m \in \Fb$ and signs $\epsilon_i \in \{\pm 1\}$ with
$R = \sum_i \epsilon_i F_i$, and to exhibit them. Since $\Sem{m+1}$ only
sees $x$-coordinates, a solution of
$\Sem{m+1}(x_1,\ldots,x_m,x(R)) = 0$ with each $x_i \in x(\Fb)$ gives a
decomposition up to signs, which are then recovered by evaluating the
group law.

\subsection{The Gr\"obner route}

Over a prime field $\FF_p$ the factor base is typically the set of points
with $x$-coordinate in an interval, and the decomposition system is
\begin{equation}
\label{eq:ideal}
  I_R = \big\langle\, \Sem{m+1}(x_1,\ldots,x_m,x_R),\;
        V(x_1),\;\ldots,\;V(x_m) \,\big\rangle,
  \qquad V(X) = \prod_{F \in \Fb}\big(X - x(F)\big).
\end{equation}
$I_R$ is zero-dimensional; $R$ is $m$-decomposable iff $I_R \neq
\langle 1 \rangle$, and a Gr\"obner basis of $I_R$ in lexicographic
order exposes the solutions. Over $\FF_{2^n}$ one instead restricts the
$x_i$ to an $\FF_2$-subspace $V$ of dimension $\ell$, writes each $x_i$
in a basis of $V$, and Weil-descends $\Sem{m+1}$ to a system of $n$
Boolean equations in $m\ell$ Boolean unknowns, again solved by F4/F5
with field equations adjoined~\cite{fppr2012,petitquisquater2012}.

In both settings the Gr\"obner solve is the step that dominates the
per-target cost, and in both settings the solve is doing real work: the
variety of $I_R$ is not given to it, it has to be found. The point of
this paper is that for small $m$ the variety is given, by the group law.

\subsection{What index calculus does and does not achieve}

Two facts frame everything that follows. Over prime fields, with
$\abs{\Fb} \approx p^{1/m}$, the probability that a random $R$
decomposes is about $1/m!$, and the number of relations needed is about
$\abs{\Fb}$; the resulting cost is $\Omega(p^{1/m} \cdot m!)$ decompositions
plus the cost of each decomposition, and Gaudry~\cite{gaudry2009} and
Diem~\cite{diem2011} showed this does not beat Pollard
rho~\cite{pollard1978} on prime-field curves for any $m$ once the
decomposition cost is accounted for. Over $\FF_{2^n}$ with $n$ prime,
the Frobenius-invariant factor base relevant to Koblitz curves has
dimension $\ell = \operatorname{ord}_2(n)$ and size $2^\ell$, which for
the cryptographic $n$ of interest is itself exponential. Faster
decomposition changes neither fact. We return to this in
\S\ref{sec:boundary}.

\section{Root-Finding Decomposition}
\label{sec:algorithms}

\subsection{Prime curves: the $\Sem{3}$ roots are $x(R \mp F_i)$}

\begin{proposition}
\label{prop:s3roots}
Let $E: y^2 = x^3 + ax + b$ over $\FF_p$, $p > 3$, let $R = (x_R, y_R)$
and $F = (x_F, y_F)$ be affine points with $x_R \neq x_F$. Then the two
roots of $X \mapsto \Sem{3}(x_R, x_F, X)$ in $\overline{\FF}_p$ are
exactly $x(R + F)$ and $x(R - F)$, counted with multiplicity.
\end{proposition}

\begin{proof}
By the defining property of $\Sem{3}$, $\Sem{3}(x_R, x_F, X) = 0$ iff
there exist $P, Q, S \in E(\overline{\FF}_p)$ with $x(P) = x_R$,
$x(Q) = x_F$, $x(S) = X$ and $P + Q + S = \mathcal{O}$. The points with
$x$-coordinate $x_R$ are $\pm R$, those with $x$-coordinate $x_F$ are
$\pm F$, so $S = -(\pm R \pm F)$ and $X \in \{x(R+F), x(R-F)\}$. The
polynomial has degree $2$ in $X$ since $x_R \neq x_F$ makes the leading
coefficient $(x_R - x_F)^2$ non-zero, so these are all the roots; if
$x(R+F) = x(R-F)$ the root is double.
\end{proof}

Proposition~\ref{prop:s3roots} says that the univariate root-finding
problem one would solve for each factor-base element in a
``root-finding'' index calculus is trivial: the roots are produced by
one slope computation. Algorithm~\ref{alg:prime} makes this concrete,
with Montgomery's batched-inversion trick~\cite{montgomery1987} so that
a sweep over the whole factor base costs a single field inversion.

\begin{algorithm}[t]
\caption{Prime-curve $2$-decomposition by direct subtraction}
\label{alg:prime}
\begin{algorithmic}[1]
\Require Curve $E: y^2 = x^3 + ax + b$ over $\FF_p$; factor base
  $\Fb = \{F_1,\ldots,F_N\}$ with $F_i = (x_i, y_i)$; hash map
  $\mathsf{idx}: x(\Fb) \to \{1,\ldots,N\}$; target $R = (x_R, y_R)$.
\Ensure $(i, \epsilon, j, \tau)$ with $R = \epsilon F_i + \tau F_j$, or
  $\bot$ if no such pair exists.
\State $D \gets [\,x_R - x_i : x_R \neq x_i\,]$ \Comment{non-zero differences}
\State $D^{-1} \gets \textsc{BatchInverse}(D)$ \Comment{one inversion, $3(N-1)$ mults}
\For{each $i$ with $x_R \neq x_i$}
  \For{$(\epsilon, \nu) \in \{(+1,\, y_R + y_i),\; (-1,\, y_R - y_i)\}$}
    \State $\lambda \gets \nu \cdot (x_R - x_i)^{-1}$ \Comment{slope of the line through $R$ and $-\epsilon F_i$}
    \State $x_3 \gets \lambda^2 - x_R - x_i$ \Comment{$= x(R - \epsilon F_i)$}
    \If{$x_3 \in \operatorname{dom}\mathsf{idx}$}
      \State $j \gets \mathsf{idx}(x_3)$;\;
             $y_3 \gets \lambda (x_R - x_3) - y_R$
      \State $\tau \gets +1$ if $y_3 = y_j$ else $-1$
      \State \Return $(i, \epsilon, j, \tau)$
    \EndIf
  \EndFor
\EndFor
\State \Return $\bot$
\end{algorithmic}
\end{algorithm}

\paragraph{Correctness and completeness.} If $R = \epsilon F_i + \tau
F_j$ with $x_i \neq x_R$, then $R - \epsilon F_i = \tau F_j$ has
$x$-coordinate $x_j \in x(\Fb)$, so the lookup at iteration $(i,\epsilon)$
succeeds and the recovered $y_3$ fixes $\tau$. Conversely every returned
tuple satisfies $R - \epsilon F_i = \tau F_j$ by construction of
$(x_3, y_3)$. The only excluded case $x_i = x_R$ means $R = \pm F_i$,
which is a $1$-decomposition and is handled by a direct lookup before
the sweep.

\paragraph{Cost.} One field inversion, $3(N-1)$ multiplications for the
batch inverse, and per factor-base element two slope evaluations
(each one multiplication, one squaring, a handful of additions) and two
hash lookups: $\Theta(N)$ field multiplications and $2N$ lookups, with
$N = \abs{\Fb}$. No polynomial arithmetic, no square roots, no linear
algebra. By contrast the Gr\"obner route computes a basis of the ideal
\eqref{eq:ideal}, whose vanishing polynomials already have degree $N$.

\paragraph{Larger $m$.} The same idea extends to $m = 3$ at cost
$\Theta(N^2)$ by sweeping pairs $(F_i, F_j)$ and looking up
$x(R - \epsilon F_i - \epsilon' F_j)$, and in general to $m$ at
$\Theta(N^{m-1})$. This is a meet-in-the-middle search rather than an
algebraic solve, and the trade-off is explicit: its cost is exponential
in $m$ with base $N$, whereas Gr\"obner methods aim at a cost polynomial
in $N$ for the structured systems that arise. In the regimes measured
here, $m = 2$, the comparison is not close. We do not claim it stays
one-sided for large $m$; see \S\ref{sec:boundary}.

\subsection{Binary curves: a meet-in-the-middle pair-sum table}

Over $\FF_{2^n}$ the factor base of interest in the Koblitz setting is a
Frobenius-invariant set $\Fb$ of size $N$, and the targets are
$m$-decomposed with $m$ up to $4$. Rather than Weil-descend $\Sem{m+1}$,
we precompute the table
\[
  T = \big\{\, x(F_i + F_j) \mapsto (i, j) \,:\, 1 \le i \le j \le N \,\big\},
\]
of size $\binom{N+1}{2}$, once per factor base. Then
(Algorithm~\ref{alg:binary}):
\begin{itemize}
\item $m = 2$: $R$ decomposes iff $x(R) \in \operatorname{dom} T$. One lookup.
\item $m = 3$: $R$ decomposes iff $x(R - F_k) \in \operatorname{dom} T$
  for some $k$. $N$ lookups.
\item $m = 4$: $R$ decomposes iff $x(R - F_k - F_l) \in \operatorname{dom} T$
  for some $k \le l$. $\binom{N+1}{2}$ lookups.
\end{itemize}
Signs are handled as in the prime case, by trying both $\pm F_k$ and
checking the recovered $y$-coordinate against the table entry. The
table is the ``left half'' of a birthday search; the right half is
enumerated per target. The procedure is exact and complete for each $m$
because it enumerates every way to write $R$ as (table entry) $+$ (at
most $m - 2$ further summands).

\begin{algorithm}[t]
\caption{Binary-curve $m$-decomposition by pair-sum table ($m \in \{2,3,4\}$)}
\label{alg:binary}
\begin{algorithmic}[1]
\Require Factor base $\Fb = \{F_1,\ldots,F_N\}$ over $\FF_{2^n}$;
  pair table $T: x(F_i + F_j) \mapsto (i,j)$ for $i \le j$; target $R$; $m$.
\Ensure Indices $(i_1,\ldots,i_m)$ and signs with $R = \sum \pm F_{i_k}$, or $\bot$.
\State $\mathcal{S} \gets \{\mathcal{O}\}$ if $m = 2$;\;
  $\{\pm F_k\}$ if $m = 3$;\;
  $\{\pm F_k \pm F_l : k \le l\}$ if $m = 4$ \Comment{right half, enumerated lazily}
\For{each $S \in \mathcal{S}$}
  \State $R' \gets R - S$
  \If{$x(R') \in \operatorname{dom} T$}
    \State $(i, j) \gets T(x(R'))$
    \State fix signs by comparing $y(R')$ with $y(F_i + F_j)$
    \State \Return $(i, j)$ together with the indices of $S$
  \EndIf
\EndFor
\State \Return $\bot$
\end{algorithmic}
\end{algorithm}

\paragraph{Cost.} Table construction is $\binom{N+1}{2}$ point
additions, paid once per factor base and amortized over all targets.
Per target the cost is $1$, $2N$ or $\sim 2N^2$ group operations and
lookups for $m = 2, 3, 4$. The matrix-F4 route pays a Macaulay-matrix
reduction per target whose size depends on the Weil-descended system,
not on $N$ directly; the measured comparison is in
\S\ref{sec:exp-binary}.

\subsection{Relationship to the Gr\"obner basis}
\label{sec:relationship}

It is worth being precise about what is and is not being replaced. The
ideal $I_R$ of \eqref{eq:ideal} is zero-dimensional and its variety
$V(I_R)$ is the set of $x$-coordinate tuples of $m$-decompositions of
$R$. A Gr\"obner basis of $I_R$ is a canonical generating set from which
$V(I_R)$ can be read off. Algorithms~\ref{alg:prime} and
\ref{alg:binary} compute $V(I_R)$ \emph{directly}, by enumeration
guided by the group law, without computing a generating set of $I_R$ at
all. Both approaches answer the same question and, on the same instance,
must agree; \S\ref{sec:verification} turns this into a check. The
Gr\"obner basis carries more information---it certifies
non-decomposability as $I_R = \langle 1 \rangle$, and the degree data it
exposes underlie the complexity analyses of
\cite{fppr2012,petitquisquater2012}---which is why we keep it available
as a baseline rather than deleting it.

\section{Implementation}
\label{sec:impl}

The algorithms were implemented in an existing Rust codebase for
elliptic-curve index calculus that already provided (i) a Koblitz-curve
index-calculus pipeline with a matrix-F4 decomposition solver over the
Weil-descended Semaev system and (ii) a generic Buchberger
implementation over multivariate polynomials with arbitrary-precision
coefficients. We added:

\begin{itemize}
\item \code{prime\_fast}: a prime-curve engine with a $64$-bit Montgomery
  field ($p < 2^{62}$), affine points, factor-base construction by
  lifting consecutive $x$-coordinates, batched inversion, and the
  direct-subtraction decomposition of Algorithm~\ref{alg:prime}.
\item \code{PairSumTable}: the table of Algorithm~\ref{alg:binary} and
  its $m \in \{2,3,4\}$ decomposition, over the codebase's packed
  $\FF_{2^n}$ representation ($n \le 63$).
\item A \code{DecompositionStrategy} enumeration in the Koblitz
  pipeline, with \code{PairTable} now the default and \code{Groebner}
  retained as a selectable baseline.
\item An arbitrary-precision (\code{BigUint}) instance of
  Algorithm~\ref{alg:prime} for fields past the native word size, and a
  P-256-specific instance over the codebase's constant-time Montgomery
  field for that prime.
\end{itemize}

The last point is the ``specialization per curve'' that the NIST
measurements call for. Each NIST prime has a Solinas form admitting a
dedicated fast reduction; the codebase already contained such a field
for P-256, and the oracle is generic over the field, so instantiating
it there was a matter of supplying the field operations. We did not
write new Solinas fields for P-192, P-224, P-384 or P-521; those curves
are measured on the generic field, and their specialization is the same
mechanical step. Nothing in the algorithm depends on the curve beyond
$a = -3$, which all five share.

\section{Fail-Closed Verification}
\label{sec:verification}

Every measurement in this paper runs under the following protocol. A
decomposition that fails any check aborts the run; no row in any table
below was produced by a run with a failing check.

\begin{enumerate}
\item \textbf{Group re-addition.} Every reported tuple
  $(i, \epsilon, j, \tau)$ is re-evaluated as $\epsilon F_i + \tau F_j$
  with an independent affine addition routine and compared to $R$ in
  both coordinates.
\item \textbf{Summation-polynomial check.} For every reported pair,
  $\Sem{3}(x_R, x_i, x_j)$ is evaluated with the codebase's existing
  $\Sem{3}$ routine and required to be zero.
\item \textbf{Variety membership.} Where a Gr\"obner basis of $I_R$ is
  also computed (the head-to-head experiments), every element of the
  reduced basis is evaluated at $(x_i, x_j)$ and required to vanish, and
  the basis is required not to be $\{1\}$.
\item \textbf{Agreement on refutation.} Where the root finder reports
  $\bot$ and a Gr\"obner basis is computed, the reduced basis is
  required to be $\{1\}$. Where the root finder reports $\bot$ on a
  target constructed as $F_i \pm F_j$, the run aborts.
\item \textbf{Positive controls.} Each experiment includes targets
  constructed as $F_i \pm F_j$ for random $i \neq j$ that must
  decompose, and random on-curve targets whose verdicts are recorded.
\end{enumerate}

Checks 3 and 4 are what make the Gr\"obner baseline more than a timer:
they establish on every target that the two methods compute the same
variety.

\section{Experiments}
\label{sec:experiments}

All timings are single-threaded wall-clock on a desktop machine, in
release builds, with the Gr\"obner baselines taken from the same
codebase as the root finders so that field arithmetic, data structures
and allocator are shared. Raw console output and JSON for every table
are in the companion repository (\S\ref{sec:repro}). Ratios are stable
to within about $2\times$ across runs; absolute times vary more on a
shared machine, and we note the one place this is visible.

\subsection{Binary Koblitz curves: pair table versus matrix-F4}
\label{sec:exp-binary}

Koblitz curves $K_0 : y^2 + xy = x^3 + 1$ over $\FF_{2^n}$ with
Frobenius-invariant factor bases; $m = 2$; all targets constructed as
$F_i + F_j$ so that every one is decomposable and both oracles must find
it. Table~\ref{tab:binary} reports per-decomposition time.

\begin{table}[h]
\centering
\caption{Binary Koblitz $m=2$ decomposition, per target. Both oracles
agreed on every target.}
\label{tab:binary}
\begin{tabular}{@{}lrrrrr@{}}
\toprule
curve & $\abs{\Fb}$ & targets & pair table & matrix-F4 & factor \\
\midrule
$K_0/\FF_{2^{9}}$  & 55   & 254 & 2.9 $\mu$s  & 1.10 ms & $383\times$ \\
$K_0/\FF_{2^{13}}$ & 4005 & 256 & 9.9 $\mu$s  & 15.3 ms & $1535\times$ \\
\bottomrule
\end{tabular}
\end{table}

\paragraph{An honest end-to-end caveat.} On a complete discrete-log
computation over $\FF_{2^9}$ with $\abs{\Fb} = 55$, the pair-table
strategy took 23.3 ms and matrix-F4 took 4.1 ms. The reason is the
one-time table build: with $\abs{\Fb} = 55$ only a few hundred
decompositions are needed, too few to amortize $\binom{56}{2}$ point
additions against a solver that is slow per target but needs no set-up.
Both strategies recovered the same logarithm. The per-decomposition
numbers in Table~\ref{tab:binary} isolate the oracle cost, which is the
quantity the Gr\"obner solve dominated; the end-to-end run is a
correctness cross-check, not a speed claim at that size.

\subsection{Small prime curves: direct subtraction versus Buchberger}
\label{sec:exp-prime-small}

Curves $y^2 = x^3 - 3x + b$ over $16$-, $20$- and $24$-bit primes with
prime group order; $\abs{\Fb} = 6$ (the Buchberger baseline is
impractical much past this); $40$ targets $[k]G$, $k = 1,\ldots,40$,
per size. The Gr\"obner baseline computes a reduced Gr\"obner basis of
$I_R$ in graded reverse-lexicographic order. Table~\ref{tab:prime-small}
reports per-target time. In each size exactly one of the $40$ targets
was decomposable and the remaining $39$ produced the unit ideal from
Buchberger, matching the root finder's $\bot$.

\begin{table}[h]
\centering
\caption{Prime $m=2$ decomposition, $\abs{\Fb}=6$, per target.}
\label{tab:prime-small}
\begin{tabular}{@{}rrrrrr@{}}
\toprule
bits & targets & decomposed / refuted & direct subtraction & Buchberger & factor \\
\midrule
16 & 40 & 1 / 39 & 0.74 $\mu$s & 26.3 ms & $35592\times$ \\
20 & 40 & 1 / 39 & 0.83 $\mu$s & 28.7 ms & $34602\times$ \\
24 & 40 & 1 / 39 & 0.65 $\mu$s & 30.6 ms & $46911\times$ \\
\bottomrule
\end{tabular}
\end{table}

The root finder was also run at factor-base sizes where the Gr\"obner
baseline is infeasible, as the decomposition oracle inside a full
index-calculus discrete-log computation, recovering the correct
logarithm at $20$ and $24$ bits with $\abs{\Fb} = 120$ (4\,438 and
78\,159 targets in 9 ms and 163 ms) and at $28$ bits with
$\abs{\Fb} = 160$ (1\,040\,671 targets in 2.83 s).

\subsection{Scaling the prime root finder}
\label{sec:exp-scaling}

To separate the method from the word size, Algorithm~\ref{alg:prime}
was run on the native $64$-bit Montgomery field at $40$, $48$, $56$ and
$61$ bits with $\abs{\Fb} = 120$, and on the arbitrary-precision field
at $p = 2^{83} - 55$ with $\abs{\Fb} = 80$. Table~\ref{tab:scaling}
reports per-target time on $200$ random on-curve targets and the
verification count on constructed positive controls.

\begin{table}[h]
\centering
\caption{Prime root finder across field sizes. ``verified'' counts
constructed $F_i + F_j$ controls that decomposed and re-added correctly.}
\label{tab:scaling}
\begin{tabular}{@{}lrrrr@{}}
\toprule
field & bits & $\abs{\Fb}$ & per target & verified \\
\midrule
native $u64$ Montgomery & 40 & 120 & 2.95 $\mu$s & 200/200 \\
                        & 48 & 120 & 3.87 $\mu$s & 200/200 \\
                        & 56 & 120 & 3.35 $\mu$s & 200/200 \\
                        & 61 & 120 & 8.82 $\mu$s & 200/200 \\
\addlinespace
arbitrary precision     & 83 & 80  & 5.99 ms     & 100/100 \\
\bottomrule
\end{tabular}
\end{table}

The native path is flat to the $2^{62}$ ceiling. The $83$-bit row is
three orders of magnitude slower because that first arbitrary-precision
instance performed one un-batched modular inversion per factor-base
element; the NIST instances below use the batched form and recover most
of the gap. The point of the row is that the algorithm, not the word
size, is what extends.

On the binary side the pairs-and-solve root finder over a bounded
subspace of $\FF_{2^{53}}$ was cross-checked against exhaustive $\Sem{4}$
search on $200/200$ targets at subspace dimension $\ell = 6$, and its
rejection cost measured at $\ell = 8, 10, 11$ (38 ms, 585 ms, 2.48 s per
target, growing with $\abs{V} = 2^\ell$ as expected). The subspace is
bounded on purpose; see \S\ref{sec:boundary}.

\subsection{NIST prime curves P-192 through P-521}
\label{sec:exp-nist}

The five NIST prime curves~\cite{fips186-4} were instantiated with their
standard $p$ and $b$ ($a = -3$ throughout); the factor base is the first
$N$ points obtained by lifting $x = 1, 2, \ldots$. Three measurements
were made per curve.

\paragraph{A. Head-to-head against Buchberger.} $\abs{\Fb} = 6$; $10$
targets per curve, $5$ constructed as $F_i + F_j$ and $5$ random
on-curve points. Both methods run in the same arbitrary-precision field.
All checks of \S\ref{sec:verification} passed on every target: all $5$
constructed targets decomposed and verified on every curve, all $5$
random targets were refuted by both methods on every curve, and the
root finder's witnesses lay in the variety of the reduced Gr\"obner
basis in every case.

\begin{table}[h]
\centering
\caption{NIST P-curves, $m=2$, $\abs{\Fb}=6$, $10$ targets each, same
field and ideal. All fail-closed checks passed.}
\label{tab:nist-f4}
\begin{tabular}{@{}lrrrrr@{}}
\toprule
curve & bits & direct subtraction & Buchberger & factor & decomposed / refuted \\
\midrule
P-192 & 192 & 120 $\mu$s & 414 ms & $3455\times$ & 5 / 5 \\
P-224 & 224 & 218 $\mu$s & 1.27 s & $5802\times$ & 5 / 5 \\
P-256 & 256 & 203 $\mu$s & 613 ms & $3021\times$ & 5 / 5 \\
P-384 & 384 & 337 $\mu$s & 1.09 s & $3235\times$ & 5 / 5 \\
P-521 & 521 & 747 $\mu$s & 2.54 s & $3395\times$ & 5 / 5 \\
\bottomrule
\end{tabular}
\end{table}

The factor is roughly flat in the field size. Both methods pay the same
arbitrary-precision arithmetic per operation; what separates them is
that the root finder performs $\Theta(\abs{\Fb})$ field operations
while the Gr\"obner solve performs a basis computation whose
coefficient arithmetic is all in that field.

\paragraph{B. Root-finder scaling at $\abs{\Fb} = 64$.} $200$ random
on-curve targets (the rejection path, which dominates in practice) and
$100$ constructed controls $F_i \pm F_j$, every control re-added in the
group and checked against $\Sem{3}$.

\begin{table}[h]
\centering
\caption{NIST P-curves, root finder only, $\abs{\Fb}=64$, generic
arbitrary-precision field.}
\label{tab:nist-scale}
\begin{tabular}{@{}lrrr@{}}
\toprule
curve & rejection / target & controls verified & control / target \\
\midrule
P-192 & 369 $\mu$s & 100/100 & 301 $\mu$s \\
P-224 & 428 $\mu$s & 100/100 & 268 $\mu$s \\
P-256 & 433 $\mu$s & 100/100 & 324 $\mu$s \\
P-384 & 583 $\mu$s & 100/100 & 394 $\mu$s \\
P-521 & 1.28 ms    & 100/100 & 3.66 ms$^{*}$ \\
\bottomrule
\end{tabular}

\smallskip
{\footnotesize $^{*}$Run-to-run jitter on a shared machine; the
preceding run measured 1.13 ms. The ratios in Table~\ref{tab:nist-f4}
are stable to within $\sim 2\times$.}
\end{table}

\paragraph{C. Per-curve specialization on P-256.} The identical oracle,
with the identical batched inversion, instantiated over the codebase's
dedicated constant-time Montgomery field for P-256.

\begin{table}[H]
\centering
\caption{P-256, $\abs{\Fb}=64$: generic field versus dedicated field.
Both paths agreed on every target.}
\label{tab:nist-special}
\begin{tabular}{@{}lrrr@{}}
\toprule
 & generic & P-256 field & factor \\
\midrule
rejection / target & 433 $\mu$s & 41 $\mu$s & $11\times$ \\
control / target   & 324 $\mu$s & 28 $\mu$s & $12\times$ \\
\bottomrule
\end{tabular}
\end{table}

\paragraph{Closing check.} The benchmark asserts that the root finder's
total time is below Buchberger's on every one of the five curves with
all verification checks passing, and exits non-zero otherwise. It
passed.

\section{What the Measurements Do and Do Not Show}
\label{sec:boundary}

We state the boundary of the result in the same place as the result.

\paragraph{It is a constant-factor improvement to a subroutine.} The
number of decompositions index calculus must attempt is unchanged, and
so is the probability that a random target decomposes. Over prime
fields the cost remains on the order of $p/\abs{\Fb}$ group operations
per relation for $m = 2$, and the Gaudry--Diem analysis that
prime-field index calculus does not beat Pollard rho for any $m$ is
untouched~\cite{gaudry2009,diem2011}. The oracle was the slow step
\emph{per target}; it was never the reason prime-field index calculus
fails to be competitive.

\paragraph{The exponential factor base on Koblitz curves is untouched.}
The Frobenius-invariant factor base relevant to $K_0/\FF_{2^n}$ has
dimension $\ell = \operatorname{ord}_2(n)$: $\ell = 52$ at $n = 53$,
$\ell = 82$ at $n = 83$, so $\abs{V} = 2^{52}$ or $2^{82}$. No
decomposition engine, Gr\"obner or root-finding, materializes such a
factor base. The binary head-to-head of \S\ref{sec:exp-binary} lives
where $\operatorname{ord}_2(n)$ is small; the $\FF_{2^{53}}$
measurements of \S\ref{sec:exp-scaling} deliberately bound the subspace
to isolate the oracle cost from this wall. Raising $n$ makes the factor
base exponentially larger; it does not make decomposition harder, and
faster decomposition does not help.

\paragraph{The NIST measurements are not a security result.} On
P-192 through P-521 only the per-target decomposition subroutine was
measured, on a factor base of $6$ or $64$ points. The relation-collection
and linear-algebra stages of index calculus remain exponentially out of
reach at those sizes, and the asymptotics are unchanged. A fast
decomposition oracle says nothing about the discrete-logarithm problem
on these curves. We make no key-recovery, sub-rho, or deployed-curve
claim, and we have written the benchmark output, the research note and
this paper to prevent the result from being read as one.

\paragraph{The meet-in-the-middle cost grows with $m$.} The direct
method is $\Theta(\abs{\Fb}^{m-1})$; the Gr\"obner method is designed
for regimes where $m$ is large enough that this is prohibitive. Our
measurements are at $m = 2$ (and $m \le 4$ for the binary table), where
the comparison is decisively in favour of root finding. We do not claim
the ordering persists at the large $m$ that asymptotic index-calculus
analyses contemplate, and at those $m$ neither method is practical for
the reasons above.

\paragraph{The Gr\"obner baselines are the codebase's own.} The binary
baseline is the codebase's matrix-F4 on the identical Weil-descended
system; the prime baseline is its Buchberger on the identical ideal. A
state-of-the-art F4/F5 implementation would be faster than our
Buchberger by a large constant. It would not be faster than $\Theta(N)$
field operations, because it has to compute a basis of an ideal whose
generators already have degree $N$; and where our baseline capped its
step budget the measured time is a lower bound, so the stated ratios
understate rather than inflate its cost.

\section{Related Work}
\label{sec:related}

Semaev~\cite{semaev2004} introduced summation polynomials and proposed
the index-calculus application; Gaudry~\cite{gaudry2009} and
Diem~\cite{diem2011} gave the first rigorous and heuristic analyses for
curves over extension fields and established the negative result for
prime fields. Faug\`ere, Perret, Petit and Renault~\cite{fppr2012} and
Petit and Quisquater~\cite{petitquisquater2012} analysed the
Weil-descended systems over $\FF_{2^n}$ and the behaviour of F4/F5 on
them; Galbraith and Gebregiyorgis~\cite{galbraithgebregiyorgis2014}
studied symmetries that reduce the system size. Joux and
Vitse~\cite{jouxvitse2013} developed decomposition over small-degree
extension fields with a variant of F4 tuned to the summation-polynomial
systems. The observation that summation-polynomial systems can be
attacked by root finding rather than by Gr\"obner bases has appeared in
the prime-field setting in the work of Amadori, Pintore and
Sala~\cite{amadoripintoresala2018}, who propose a different polynomial
model; our contribution is narrower and more concrete---for $m = 2$ the
roots are given by the group law, so no polynomial model is needed at
all---and is accompanied by an implementation, a fail-closed protocol,
and measurements on standard curves. Montgomery's batched-inversion
trick~\cite{montgomery1987} is what makes the sweep cost one inversion.
Our stance on claim hygiene follows the tradition of
Kosters--Yeo~\cite{kostersyeo2015} and others who have insisted that
improvements to decomposition be reported together with their effect,
or lack of effect, on the overall complexity.

\section{Reproducibility}
\label{sec:repro}

All code and raw measurements are in the companion repository under the
following paths.

\begin{description}
\item[Algorithms.]
  \path{src/cryptanalysis/prime_fast.rs}
  (Algorithm~\ref{alg:prime}, native field);
  \path{src/cryptanalysis/koblitz_index_calculus.rs}
  (\code{PairSumTable}, Algorithm~\ref{alg:binary};
  \code{DecompositionStrategy}).
\item[Benchmarks.] \hfill\\
  \path{examples/semaev_root_vs_groebner_bench.rs}
  (\S\ref{sec:exp-binary}, \S\ref{sec:exp-prime-small}) \\
  \path{examples/semaev_rootfind_scaling_bench.rs}
  (\S\ref{sec:exp-scaling}) \\
  \path{examples/semaev_rootfind_nist_bench.rs}
  (\S\ref{sec:exp-nist})
\item[Raw evidence] (console transcripts and JSON). \hfill\\
  \path{research/semaev_root_vs_groebner_20261012/} \\
  \path{research/semaev_rootfind_scaling_20261012/} \\
  \path{research/semaev_rootfind_nist_20261012/}
\item[Research note] with the same tables.
  \path{RESEARCH_SEMAEV_ROOTFIND_VS_GROEBNER.md}.
\end{description}

\noindent To reproduce:
\begin{verbatim}
cargo run --release --example semaev_root_vs_groebner_bench
cargo run --release --example semaev_rootfind_scaling_bench
F4_CONSTRUCTED=5 F4_RANDOM=5 \
  cargo run --release --example semaev_rootfind_nist_bench
\end{verbatim}
The NIST benchmark exposes its knobs as environment variables
(\code{F4\_FB}, \code{F4\_CONSTRUCTED}, \code{F4\_RANDOM},
\code{SCALE\_FB}, \code{SCALE\_RANDOM}) and prints them at the top of
its output. The library test suite (\code{cargo test --release --lib})
passed with 2\,247 tests at the commit the measurements were taken from.

\section{Conclusion}

For the decomposition sizes at which elliptic-curve index calculus is
actually run in practice, the Gr\"obner-basis solve is not the right
tool for the decomposition oracle: the variety it would compute is
already parametrized by the group law. Replacing the solve with root
finding---direct subtraction over prime fields, a pair-sum table over
binary fields---is exact, complete, cheap to verify, and three to four
orders of magnitude faster per target on every instance we measured,
including the five NIST prime curves at their standard parameters. We
have made it the default in our codebase and kept the Gr\"obner solvers
as baselines. We have also tried to say clearly what this is not: it is
a better constant in a subroutine, not a change to what index calculus
can reach.

\begin{thebibliography}{99}

\bibitem{semaev2004}
I.~Semaev.
Summation polynomials and the discrete logarithm problem on elliptic curves.
IACR Cryptology ePrint Archive, Report 2004/031, 2004.

\bibitem{gaudry2009}
P.~Gaudry.
Index calculus for abelian varieties of small dimension and the elliptic curve discrete logarithm problem.
\emph{Journal of Symbolic Computation}, 44(12):1690--1702, 2009.

\bibitem{diem2011}
C.~Diem.
On the discrete logarithm problem in elliptic curves.
\emph{Compositio Mathematica}, 147(1):75--104, 2011.

\bibitem{faugere1999}
J.-C.~Faug\`ere.
A new efficient algorithm for computing Gr\"obner bases ($F_4$).
\emph{Journal of Pure and Applied Algebra}, 139(1--3):61--88, 1999.

\bibitem{faugere2002}
J.-C.~Faug\`ere.
A new efficient algorithm for computing Gr\"obner bases without reduction to zero ($F_5$).
In \emph{ISSAC 2002}, pages 75--83. ACM, 2002.

\bibitem{fppr2012}
J.-C.~Faug\`ere, L.~Perret, C.~Petit, and G.~Renault.
Improving the complexity of index calculus algorithms in elliptic curves over binary fields.
In \emph{EUROCRYPT 2012}, LNCS 7237, pages 27--44. Springer, 2012.

\bibitem{petitquisquater2012}
C.~Petit and J.-J.~Quisquater.
On polynomial systems arising from a Weil descent.
In \emph{ASIACRYPT 2012}, LNCS 7658, pages 451--466. Springer, 2012.

\bibitem{galbraithgebregiyorgis2014}
S.~D.~Galbraith and S.~W.~Gebregiyorgis.
Summation polynomial algorithms for elliptic curves in characteristic two.
In \emph{INDOCRYPT 2014}, LNCS 8885, pages 409--427. Springer, 2014.

\bibitem{jouxvitse2013}
A.~Joux and V.~Vitse.
Elliptic curve discrete logarithm problem over small degree extension fields.
\emph{Journal of Cryptology}, 26(1):119--143, 2013.

\bibitem{amadoripintoresala2018}
A.~Amadori, F.~Pintore, and M.~Sala.
On the discrete logarithm problem for prime-field elliptic curves.
\emph{Finite Fields and Their Applications}, 51:168--182, 2018.

\bibitem{kostersyeo2015}
M.~Kosters and S.~L.~Yeo.
Notes on summation polynomials.
arXiv:1503.08001, 2015.

\bibitem{pollard1978}
J.~M.~Pollard.
Monte Carlo methods for index computation (mod $p$).
\emph{Mathematics of Computation}, 32(143):918--924, 1978.

\bibitem{montgomery1987}
P.~L.~Montgomery.
Speeding the Pollard and elliptic curve methods of factorization.
\emph{Mathematics of Computation}, 48(177):243--264, 1987.

\bibitem{fips186-4}
National Institute of Standards and Technology.
Digital Signature Standard (DSS).
FIPS PUB 186-4, July 2013.

\end{thebibliography}

\end{document}
